\scanpage{561}
\mspage{285}

\[
(53)\qquad
\widetilde M_{1,1}\simeq M_{1,1}[2]'/\Gamma'_2
\simeq U_{0,3}/\mathfrak S_3
\]

\unsure{one|now}\qquad
(valid over $\mathrm{Spec}\,\mathbb Z[1/2]$)

one does indeed find, over \inserted{$S=$}$\mathrm{Spec}\,\mathbb Z[1/2]$, that
$\widetilde M_{1,1}$ is obtained from a relative scheme \inserted{$Y$,} which is
\deleted{\illegible{1}} locally \inserted{(étale)} isomorphic to $\mathbb P^1_S$,
by removing a section $s$ (whose existence already implies that $Y$ is
\unsure{globally|globally} isomorphic to $\mathbb P^1_S$ - hence
$Y-s(S)$ isomorphic to $\mathbb E^1_S$.
In this approach one thus feels like denoting by $\infty$
(not by $0$, as in the theory of maps!) this section,
which corresponds to the ``points at infinity'' $\unsure{(0,1,\infty)}$ of
$U_{0,3}$, or again to the ``\unsure{point} at infinity'' of
$\widetilde M_{1,1}$.
It is moreover easy to check a priori \unsure{that}
(by the Deligne-Mumford compactification of $M_{1,1}$ and of
$\widetilde M_{1,1}$) that $\widetilde M_{1,1}$ over the whole of $\mathrm{Spec}\,\mathbb Z$
is of the form $Y-\mathrm{Im}\,s$, where $Y$ is smooth proper
over $\mathrm{Spec}\,\mathbb Z$, and $s$ a section - hence what one
knows for example about the generic fiber, and the fact that $\mathbb Z$ is
principal, imply that $Y\simeq\mathbb P^1_{\mathbb Z}$, and

\scanpage{562}

that one therefore has $\widetilde M_{1,1}\simeq\mathbb E^1_{\mathbb Z}$.
To \deleted{it} choose this isomorphism, one uses the two sections of
$\mathbb E^1_{\mathbb Z}$, corresponding to the anharmonic curves.
These two sections \uline{are not} \uline{disjoint}, they
meet in characteristic $2$ and characteristic $3$.
(As for what happens in characteristic $\ne 2$ i.e. over
$\mathrm{Spec}\,\mathbb Z[1/2]$, one sees it clearly from the isomorphism $53$ -
\deleted{the} for in characteristic $3$ the two orbits
\[
(2,-1,1/2)\quad\text{and}\quad(j,\bar j)
\]
coincide and collapse into one and the
same point).
But taking one of these sections as zero section
(unless I am mistaken it is the one corresponding to the elliptic curve richest
in symmetries, \unsure{namely|i.e.} \unsure{of|the} group $\mathbb Z/6\mathbb Z$, that one takes -
it is a matter of convention, of course - it is therefore also the orbit
$\{j,\bar j\}$, corresponding to a stabilizer
$\simeq\mathbb Z/3\mathbb Z$) - this determines the isomorphism
$\widetilde M_{1,1}\simeq\mathbb E^1_{\mathbb Z}$ up to sign.
The other section then becomes an integer of the form
$\pm 2^a3^b$ ($a,b\in\mathbb N^*$ - \uline{integers} well determined whose
value \unsure{I|I have} do not know by \unsure{heart}), and one finishes

\scanpage{563}
\mspage{286}

normalizing by requiring that the \uline{sign} be \uline{plus}. Thus

\uline{Theorem} (for the record - this is well known) There is a unique isomorphism
\[
(54)\qquad
\widetilde M_{1,1}\simeq\mathbb E^1_{\mathbb Z},
\]
which \deleted{gives} in characteristic $0$ gives to the anharmonic curve with
automorphism group $\mathbb Z/6\mathbb Z$
the \deleted{\illegible{1}}\inserted{invariant} $0$, and to the one with
group $\mathbb Z/4\mathbb Z$ an invariant \deleted{\illegible{1}} $>0$
(which will necessarily be $2^a3^b$, with $a,b\in\mathbb N^*$ well
determined, but forgotten by me\ldots).
\margin{$b=3$, cf. below}

In \unsure{terms} of the isomorphism
\[
\deleted{\unsure{(55)}}\qquad
\widetilde M_{1,1}\simeq U_{0,3}/\mathfrak S_3
\subset \mathbb P^1/\mathfrak S_3
\]
(valid over $\mathrm{Spec}\,\mathbb Z[1/2]$)
this corresponds to \deleted{the isomorphism} an isomorphism between
$\mathbb P^1/\mathfrak S_3$ and $\mathbb P^1$ which to the orbit
$(0,1,\infty)$ assigns $\infty$ (not $0$) and to the orbit $(j,\bar j)$
assigns $0$ (not $\infty$), to the orbit $(2,-1,1/2)$
\unsure{the famous} $2^a3^b$ (and not $1$). The function
\margin{\uline{modular} \uline{invariant}}
\[
(55)\qquad j(\lambda)\in\mathbb Q(\lambda)
\]
which realizes this isomorphism, i.e. \unsure{the|the} \deleted{\unsure{morphism}} morphism
\[
(56)\qquad
\mathbb P^1_{\mathbb Z}\xrightarrow{j}\mathbb P^1_{\mathbb Z}
\]
corresponding (of degree $6$, with
\[
j\circ u=j
\qquad(u\in\mathfrak S_3)
\]
)

\scanpage{564}

is therefore \unsure{related|related} to the one inspired by the theory of maps,
i.e. $f(\lambda)$, by
\[
(57)\qquad
j(\lambda)=2^a3^b f(\lambda)^{-1}.
\]

Moreover $f(\lambda)$ \inserted{(or $j(\lambda)$)} is readily computed
from knowledge of its zeros and its poles,
\inserted{with their multiplicities},
\deleted{\illegible{2}} and the ``normalization'' by the value of $f$
(or of $j$) on one of the sections $2,-1,1/2$; one finds immediately:
\[
(58)\qquad
f(z)=\frac{3^3}{2^2}\frac{z^2(z-1)^2}{(z^2-z+1)^3}
\]
hence
\[
(59)\qquad
j(\lambda)=2^{a+2}3^{b-3}\frac{(z^2-z+1)^3}{z^2(z-1)^2}.
\]

Since $j(\lambda)$ must still make sense in characteristic $3$ and not be
constant, this shows moreover \uline{$b=3$}, so $(59)$ can also be written
\[
\deleted{\unsure{(60)}}\qquad
j(\lambda)=2^{a+2}\frac{(z^2-z+1)^3}{z^2(z-1)^2}
\]
- but one \unsure{will recall|recalls} that this formula is not relevant -
\unsure{does not allow} one to compute
\deleted{a \unsure{discriminant} of}\inserted{the invariant of an}
elliptic curve - except in characteristic $\ne 2$, fortunately!
To \deleted{\illegible{1}}

\scanpage{565}
\mspage{287}

\deleted{define}\inserted{describe}
\[
(60)\qquad
M_{1,1}\longrightarrow\mathbb E^1_{\mathbb Z}
\]
\margin{and in particular to determine
the other critical modular invariant
$2^a3^3$, i.e. \uline{determine $a$}}
also in a neighborhood of $2$ - say in characteristic $\ne 3$ - one needs a study
of elliptic curves which no longer goes by way of the representation
(of \uline{Legendre}, unless I am mistaken)
$y^{2}=\sqrt{x(x-1)(x-\lambda)}$, or again, by way of the Jacobi rigidification
of level $2$, but on the other hand by the one of level $3$.

\uline{Remarks} One sees at once that the quotient
$\mathbb P^1_{\mathbb Z}/\mathfrak S_3$ identifies with
$\mathbb P^1_{\mathbb Z}$, uniquely by taking as images of the
orbits $(0,1,\infty)$ and $(j,\bar j)$ the sections $0$ and $\infty$
(for example, following the convention corresponding to the theory of maps),
and by taking as image of the orbit $(2,-1,1/2)$ a
\deleted{integer}\inserted{rational} which is $>0$.
This is possible a priori, for one notes that the \deleted{sections}
orbits $(0,1,\infty)$ and $(j,\bar j)$
\deleted{on the one hand, $(0,1,\infty)$ and $(2,-1,1/2)$ on the other hand,}
do not coincide in any characteristic, hence correspond to
disjoint sections of $Y=\mathbb P^1_{\mathbb Z}/\mathfrak S_3$,
whereas the \unsure{remaining|constant} orbit

\scanpage{566}

$(2,-1,1/2)$ coincides with the first in characteristic $2$, with the second in
characteristic $3$ (and these are the only coincidences that can occur).
One therefore sees a priori that the third section will be of the form
$2^\alpha/3^\beta$, with $\alpha,\beta\in\mathbb N^*$.
The explicit calculation is moreover obvious, the composite morphism
\[
(61)\qquad
\mathbb P^1_{\mathbb Z}
\xrightarrow{f_1^{\unsure{*|\circ}}}
\mathbb P^1_{\mathbb Z}/\mathfrak S_3
\simeq\mathbb P^1_{\mathbb Z}
\]
must be of the form
$f_1^{\unsure{*|\circ}}=c.f^{\unsure{*|\circ}}$, with the constant
\inserted{$c\in\mathbb Q$} chosen in such a way that
$cf^{\unsure{*|\circ}}$ \deleted{has} reduces properly in
\unsure{all} characteristics (which determines $c$ \unsure{modulo|modulo the} sign)
and that
\inserted{$f_1(-1)=$}\deleted{\illegible{1}}
$cf^{\unsure{*|\circ}}(-1)$ \inserted{$=c$} $>0$
(which removes the indeterminacy of the sign).
One then finds
\[
(62)\qquad
f_1(\lambda)=
\frac{\lambda^2(\lambda-1)^2}{(\lambda^2-\lambda+1)^{2}}
=\frac{2^2}{3^3}f(\lambda)
\]
hence
\[
(63)\qquad
f_1(-1)=\frac{2^2}{3^3}
\qquad(\text{so }\alpha=2,\ \beta=3).
\]

But one will take note \deleted{that}
\inserted{that in a neighborhood of characteristic $2$,}
these calculations no longer concern $\widetilde M_{1,1}$ and
$\unsure{\widetilde{\widetilde M}|\widetilde{\widehat M}}_{1,1}$,
where the configuration of the three sections
\deleted{\illegible{1}
$\unsure{\Delta_\infty,\Delta_0|\Lambda_\infty,\Lambda_0}$,}
is not the same \unsure{as here}\ldots

\scanpage{567}
\mspage{288}

Following the \unsure{momentum} of these reflections, it would be natural to look
more closely at the scheme
\[
(64)\qquad
M_{1,1}[4],\quad
\text{on which acts }\Gamma_4=\mathrm{Gl}(2,\mathbb Z/4)
\]
\[
\text{via}\quad \Gamma'_4=\Gamma_4/\pm1
\]
and \unsure{the homomorphism} $\otimes$
\[
(65)\qquad
M_{1,1}[4]\longrightarrow M_{1,1}[2]'\simeq U_{0,3}
\]
\inserted{Compatible} \deleted{\illegible{1}}
\[
\Gamma'_4=\mathrm{Gl}(2,\mathbb Z/4)/\pm1
\longrightarrow
\Gamma'_2=\Gamma_2\simeq\mathrm{Gl}(2,\mathbb Z/2)
\simeq\mathfrak S_3
\]
\margin{\diagram{[Above the arrow, a crossed-out inscription:
\deleted{\illegible{1} $\det$}. To the left of (65), a reference line
connects the sign $\otimes$ to the marginal note.]}}
\margin{in fact the kernel of the homomorphism
$\mathrm{Gl}(k,A)$\deleted{\illegible{1}}$\longrightarrow\mathrm{Gl}(k,A/J)$
is always a commutative group $\circledast$
\unsure{if} $J$ is a square-zero ideal of a ring $A$ -
so here $\mathrm{Ker}(\mathrm{Gl}(2,\mathbb Z/4\mathbb Z)\longrightarrow
\mathrm{Gl}(2,\mathbb Z/2\mathbb Z))$ is a \uline{commutative} group,
and even a vector space over $\mathbb F_2$ - it is isomorphic to
$\mathbb F_2^4$}
\margin{$\circledast$ and even an $A/J$ $(A/J)^4$}

An immediate calculation, taking into account that
\[
(\mathbb Z/4\mathbb Z)^*\simeq\{\pm1\}
\]
, gives
\[
\mathrm{card}\,\Gamma_4
=
2\underbrace{\mathrm{card}\,\mathrm{Sl}(2,\mathbb Z/4\mathbb Z)}
_{8\,\mathrm{card}\,\mathrm{Sl}(2,\mathbb Z/2\mathbb Z)}
=
16.\underbrace{\mathrm{card}\,\Gamma_2}_{6}
\]
hence
\[
(66)\qquad
\mathrm{card}\bigl(\Gamma'_{4,2}
=\mathrm{Ker}(\Gamma'_4\longrightarrow\Gamma'_2)\bigr)=8
\]
i.e. $M_{1,1}[4]$ is a Galois covering of $M_{1,1}[2]'$,
of \deleted{group} order $8$.
Let us note that it is not geometrically connected over $\mathbb Q$,
\inserted{\unsure{or over $\mathrm{Spec}\,\mathbb Z[1/2]$}},
- since $M_{1,1}[4]$ lies over the quadratic extension
$\mu_4^*$ \deleted{\unsure{$\simeq\mathrm{Spec}$}}, of
$\mathrm{Spec}\,\underbrace{\mathbb Z[1/2]}_{\Lambda}$ say,
which is none other than $\Lambda[i]$.
\inserted{One has \deleted{\illegible{2}}\inserted{a}}
\deleted{\unsure{Let us choose the point} \illegible{1}}
\inserted{factorization of
$M_{1,1}[4]\longrightarrow M_{1,1}[2]'$}
\deleted{\unsure{over}
$\mathrm{Spec}\,\Lambda[i]=\mu_4^*$
over $\mathrm{Spec}\,\mathbb Z[1/2]$,
we can distinguish a \unsure{connected component}}

\scanpage{568}

\((67)\)
\diagram{
\begin{tikzcd}[column sep=large,row sep=large]
M_{1,1}[4]
  \arrow[r,"\text{Galois}"]
  \arrow[d]
&
(U_{0,3})_{\unsure{S'}}
  \arrow[r,"\text{Galois }\unsure{\text{degree }2}"]
  \arrow[dl]
&
M_{1,1}[2]'\simeq U_{0,3}
  \arrow[d]
\\
S'=\mu_4^*
  \arrow[rr,"\text{Galois}","\unsure{\text{degree }2}"']
&&
\mathrm{Spec}\,\mathbb Z[1/2]=S
\end{tikzcd}
}

where now
\[
M_{1,1}[4]\longrightarrow M_{1,1}[2]'_{\unsure{S'}}
=(U_{0,3})_{S'}
\]
is a Galois covering with group the kernel of
$\Gamma'_{4,2}\xrightarrow{\det}\{\pm1\}$,
\deleted{\illegible{1}}\inserted{or again} that of
$\mathrm{Sl}(2,\mathbb Z/4\mathbb Z)'
\longrightarrow\mathrm{Sl}(2,\mathbb Z/2\mathbb Z)'$,
which we shall denote by $S\Gamma'_{4,2}$, which is of order $4$.

\uline{Proposition} The group
\[
(68)\qquad
S\Gamma_{4,2}
=
\mathrm{Ker}\bigl(
\mathrm{Sl}(2,\mathbb Z/4\mathbb Z)
\longrightarrow
\mathrm{Sl}(2,\mathbb Z/2\mathbb Z)
\bigr)
\]
is isomorphic to \deleted{\illegible{2}}
$\unsure{\mathfrak{sl}}(2,\mathbb F_2)$
(group of matrices of trace zero with coefficients in $\mathbb F_2$),
\inserted{with \unsure{the obvious action|the obvious actions}
of $\mathrm{Sl}(2,\mathbb F_2)$ on it,}
and the subgroup $\pm1$ corresponds to scalar matrices -
the quotient \deleted{\unsure{is}}\inserted{$S\Gamma'_{4,2}$ of the two}
isomorphic to \deleted{\illegible{1}} $\mathbb F_2^2$
(in a silly way and not canonical at all!).
\margin{\uline{N.B.} It is the $4$ complements of this \unsure{subspace}
$\mathbb F_2\unsure{\hookrightarrow|\subset}\unsure{\mathfrak{sl}}(2,\mathbb F_2)$
\deleted{(\unsure{commutative}}
which form a torsor \unsure{under|over} the dual of
$V=\unsure{\mathfrak{sl}/\mathbb F_2|\mathfrak{sl}(\mathbb F_2)}$
which should \unsure{correspond to the four} \deleted{\illegible{1}}
liftings of $\pi_{0,3}\hookrightarrow T^{+\prime}_{1,1}=\mathrm{Sl}(2,\mathbb Z)'$,
\unsure{i.e.|to} $\pi_{0,3}\hookrightarrow\mathrm{Sl}(2,\mathbb Z)$
cf. pages $259$, $260$ \unsure{above} on these questions}
\deleted{\unsure{finally}}
Viewing $M_{1,1}[4]$ as a Galois covering of
$M_{1,1}[2]'_{S'}=(U_{0,3})_{S'}$, one finds
\uline{geometrically} ``the'' universal \uline{abelian} covering of
$U_{0,3}$ with group killed by $2$ (and for good reason), which has as
Galois group $(\mathbb F_2)^{\{0,1,\infty\}}/\mathrm{Diagonale}$.

\scanpage{569}
\mspage{289}

Finally, everything is essentially tautological - but an amusing calculation to
do will be to ``\inserted{re}describe the action of $\mathrm{Sl}(2,\mathbb F_2)$ on
$\unsure{\mathfrak{sl}}(2,\mathbb F_2)/\text{scalar matrices}$''
in such a way as to identify it with the action of $\mathfrak S_3$ on
$\mathbb F_2^{\{0,1,\infty\}}/\text{diagonal}$.

As a covering of \deleted{$\unsure{\widetilde M}_{1,1}\simeq M_{1,1}[4]$}
\[
(69)\qquad
\widetilde M_{1,1\,S'}
\simeq M_{1,1}[4]/S\Gamma'_4
\simeq M_{1,1}[2]'_{\unsure{S'}}/\Gamma'_2,
\]
$M_{1,1}[4]$ is a Galois covering of order $24$ ($=4.6$),
with group $\mathrm{Sl}(2,\mathbb Z/4\mathbb Z)'$.
\deleted{It is} It would be quite \unsure{rare|welcome} if one did not have an isomorphism
\[
(70)\qquad
\mathrm{Sl}(2,\mathbb Z/4\mathbb Z)'
\simeq\mathfrak S_4
\qquad
(\text{automorphisms of the octahedron})
\]
and if $M_{1,1}[4]_{S'}$ were its covering \deleted{\unsure{étale}}
\deleted{\illegible{1}}
of $\mathbb P^1_{S'}\setminus\{0,1,\infty\}$, with
ramification compatible with
$2\{1\}$ \deleted{\illegible{1}} $+3\{\infty\}$,
which corresponds to the \uline{octahedron} \unsure{\&} in the theory of
maps - one should have over
$S'=\mathrm{Spec}\,\mathbb Z[1/2][i]$
a \inserted{canonical} isomorphism
\deleted{\unsure{of $\widetilde M_{1,1}$}}
\[
(71)\qquad
M_{1,1}[4]_{S'}\simeq
\text{Fermat curve}
\quad
x^2+y^2+z^2=0
\quad(\text{over }S')
\]
\margin{it would be better to write the Fermat equation
here
$x_0^2+x_1^2+x_\infty^2=0$}
(whose automorphism group is \unsure{precisely} an extension of
$\mathfrak S_3$ by
$\mu_2^{\{0,1,\infty\}}/\mathrm{diag}$!).

\scanpage{570}

To do things properly, one would have to make explicit the relation between elliptic
curves \inserted{$E/S$} (over a scheme $S$ over
$\mathbb Z[1/2][i]$) \unsure{equipped} with a Jacobi rigidification
of level $4$ $e_1,e_2$ with $e_1\wedge e_2=[i]^{-1}$, and
``nontrivial'' solutions of the equation
\[
(72)\qquad x^2+y^2+z^2=0
\]
over $S$ (i.e. the systems of sections
$x,y,z\in\Gamma(S,\mathcal O_S^*)$ satisfying $(72)$, modulo
multiplication by a ``scalar''
\inserted{$\alpha\in\Gamma(S,\mathcal O_S^*)$}
- or again, breaking the symmetry by putting
$X=x/z$, $Y=y/z$, \deleted{\illegible{1}}
the solutions of the equation
\[
(73)\qquad X^2+Y^2=-1\qquad(=i^2)
\]
$(X,Y\in\Gamma(S,\mathcal O_S^*))$\ldots

One has thus interpreted $M_{1,1}[2]'$, and
(modulo some checks and a somewhat more thorough study)
$M_{1,1}[4]$, in relation with \deleted{two} regular
\inserted{weighted triangular}
maps - the ``\unsure{dihedral|dihedral}'' map
\inserted{\unsure{of index $3$|of order $3$}}
or universal weighted triangular map
(three vertices $0,1,\infty$,
\figure{[An almost circular outline, with a horizontal oval.
Three filled dots are marked on the oval: $0$ on the left,
$\infty$ on the upper branch on the right, $1$ on the
lower branch on the right. The three portions of the oval join these
three points.]}
three edges, two faces which are triangles, type
$(p,q)=(3,3)$), and the \uline{octagonal} map,
which is a covering of order $4$ of this one
\unsure{(}\deleted{\illegible{2}}\inserted{étale}
away from the vertices $\{0,1,\infty\}$),

There would be \unsure{reason} also to look at